\chapter{Unitary circuits}
\label{ch:unitary_circuits}

This chapter uses the notation of Chapters~\ref{ch:sites} and~\ref{ch:circuits}:
$V$ is a finite set of sites, $X=(X_v)_{v\in V}$ a system of finite basis sets,
$\mathcal H_X=\bigotimes_v\C^{X_v}$, and $\Gamma$ a geometry on $V$. A unitary
circuit is a circuit in the sense of Chapter~\ref{ch:circuits} in which every
layer acts on one system $X$ and every gate and every onsite family consists of
a single unitary. Local circuits of depth $T$ on a ring
\cite{Malz2024Preparation}, depth-$\ell$ circuits on the edges of a lattice
\cite{Piroli2021LOCC} and products of two-site gates on an open chain are such
circuits on three geometries, measured by depth or by gate count. A one-site
unitary is an onsite family and costs nothing. The identity, composition,
monotonicity and the light cone are the statements of
Chapter~\ref{ch:circuits} for these circuits.

\section{Graphs and windows}
\label{sec:unitary_graph_geometry}

\begin{definition}[Graph and window geometries]
    \label{def:unitary_graph_geometry}
    \lean{QuantumCircuit.Geometry.graph, QuantumCircuit.Geometry.mem_graph_cell,
        QuantumCircuit.Geometry.graph_cell_mk, QuantumCircuit.Geometry.graph_cost_pos,
        QuantumCircuit.Geometry.pathCell, QuantumCircuit.Geometry.pathCell_cell,
        QuantumCircuit.Geometry.exists_pathCell, QuantumCircuit.Geometry.window,
        QuantumCircuit.Geometry.mem_window_cell}
    \leanok
    \uses{def:circuits_geometry}
    For a simple graph $G$ on $V$, the graph geometry has one cell $\{u,v\}$ of
    cost $1$ for each edge $uv$ \cite{Piroli2021LOCC}. The open chain of $n$ sites
    is the graph geometry of the path on $\{0,\dots,n-1\}$, whose cells are the
    pairs $\{p,p+1\}$. The windows of $m$ consecutive sites of this chain are the
    cells $\{p: a\le p<a+m\}$ of cost $1$, one for each $a<n$, cut off at the
    right end of the chain.
\end{definition}

\begin{lemma}[Light cones of a graph and of the open chain]
    \label{lem:unitary_graph_light_cone}
    \lean{QuantumCircuit.Geometry.iterate_neighbourhood_graph_subset,
        QuantumCircuit.Geometry.disjoint_iterate_neighbourhood_graph,
        QuantumCircuit.Geometry.iterate_neighbourhood_graph_subset_dist,
        QuantumCircuit.Geometry.disjoint_iterate_neighbourhood_graph_of_dist,
        QuantumCircuit.Geometry.iterate_neighbourhood_pathGraph_subset,
        QuantumCircuit.Geometry.disjoint_iterate_neighbourhood_pathGraph}
    \leanok
    \uses{def:unitary_graph_geometry, lem:circuits_neighbourhood_mono,
        lem:sites_path_distance}
    In the graph geometry of $G$, the $r$-fold neighbourhood $\mathcal N^r(R)$ of a
    region $R$ lies within graph distance $r$ of $R$, and regions at graph distance
    larger than $2r$ have disjoint $r$-fold neighbourhoods. On the open chain the
    same holds with the distance $|i-j|$ of positions.
\end{lemma}
\begin{proof}\leanok
    By induction on $r$: a site of $\mathcal N^{r+1}(R)$ lies in $\mathcal N^r(R)$ or in
    an edge meeting a site $u$ of $\mathcal N^r(R)$, and the two ends of an edge are at
    distance at most $1$; the triangle inequality gives the bound. If $v$ lay in both
    neighbourhoods, the triangle inequality would join $R$ and $R'$ within $2r$. On the
    path the graph distance is the distance of positions.
\end{proof}

\section{Unitary circuits and their costs}
\label{sec:unitary_circuits}

\begin{definition}[Unitary circuit]
    \label{def:unitary_circuit}
    \lean{QuantumCircuit.UnitaryCircuit, QuantumCircuit.UnitaryCircuit.toCircuit,
        QuantumCircuit.UnitaryCircuit.op, QuantumCircuit.UnitaryCircuit.depth,
        QuantumCircuit.UnitaryCircuit.gateCount,
        QuantumCircuit.UnitaryCircuit.instUniqueIndex,
        QuantumCircuit.UnitaryCircuit.kraus_eq_op,
        QuantumCircuit.UnitaryCircuit.forall_isUnitary,
        QuantumCircuit.UnitaryCircuit.forall_isTP,
        QuantumCircuit.UnitaryCircuit.isUnbranched,
        QuantumCircuit.Layer.nonempty_unique_index_of_forall_isUnitary}
    \leanok
    \uses{def:circuits_circuit, def:circuits_local, def:circuits_costs,
        lem:circuits_seq, lem:circuits_unique_branch}
    A unitary circuit on $\Gamma$ on the system $X$ is a finite sequence of layers
    from $X$ to $X$ whose gates and onsite families are each a single unitary
    \cite{Malz2024Preparation, Piroli2021LOCC}. As a circuit it is the sequence of
    its layers; it has exactly one branch, and its operator $U$ is the Kraus
    operator of that branch, $U=W_T\cdots W_1$ for the operators $W_k$ of its
    layers. Its depth and gate count are those of the circuit.
\end{definition}

\begin{lemma}[Costs of a unitary circuit]
    \label{lem:unitary_circuit_costs}
    \lean{QuantumCircuit.UnitaryCircuit.nonempty_index,
        QuantumCircuit.UnitaryCircuit.depth_eq_sum,
        QuantumCircuit.UnitaryCircuit.gateCount_eq_sum,
        QuantumCircuit.UnitaryCircuit.depth_le_gateCount,
        QuantumCircuit.UnitaryCircuit.dual_eq,
        QuantumCircuit.UnitaryCircuit.op_mem_unitaryGroup,
        QuantumCircuit.UnitaryCircuit.conjTranspose_op_mul_op,
        QuantumCircuit.UnitaryCircuit.op_mul_conjTranspose_op}
    \leanok
    \uses{def:unitary_circuit, thm:circuits_cost_subadditive, thm:circuits_tp}
    The depth and the gate count of a unitary circuit are the sums of those of its
    layers, and the depth is at most the gate count. Its Heisenberg picture is
    $A\mapsto U^\dagger AU$, and $U$ is unitary.
\end{lemma}
\begin{proof}\leanok
    Every layer has a Kraus index, so the cost of the sequence is the sum of the
    costs of the layers. The circuit has one branch, so its Heisenberg picture is
    conjugation by $U$; its families are trace-preserving, so the Heisenberg picture
    fixes $1$, which is $U^\dagger U=1$, and $UU^\dagger=1$ follows for square
    matrices.
\end{proof}

\begin{theorem}[Identity, composition and inverse]
    \label{thm:unitary_circuit_composition}
    \lean{QuantumCircuit.UnitaryCircuit.id, QuantumCircuit.UnitaryCircuit.op_id,
        QuantumCircuit.UnitaryCircuit.depth_id, QuantumCircuit.UnitaryCircuit.gateCount_id,
        QuantumCircuit.UnitaryCircuit.comp, QuantumCircuit.Circuit.seq_append,
        QuantumCircuit.UnitaryCircuit.toCircuit_comp, QuantumCircuit.UnitaryCircuit.op_comp,
        QuantumCircuit.UnitaryCircuit.depth_comp,
        QuantumCircuit.UnitaryCircuit.gateCount_comp,
        QuantumCircuit.UnitaryCircuit.adjoint, QuantumCircuit.UnitaryCircuit.op_adjoint,
        QuantumCircuit.UnitaryCircuit.depth_adjoint,
        QuantumCircuit.UnitaryCircuit.gateCount_adjoint,
        QuantumCircuit.Layer.forall_isUnitary_adjoint, QuantumCircuit.Layer.depth_adjoint,
        QuantumCircuit.Layer.gateCount_adjoint}
    \leanok
    \uses{def:unitary_circuit, lem:unitary_circuit_costs, thm:circuits_composition,
        lem:circuits_adjoint}
    The empty unitary circuit has operator $1$ and cost $0$. The circuit $c$
    followed by $c'$ has operator $U'U$, depth the sum of the depths and gate count
    the sum of the gate counts. The adjoint layers of $c$ in reverse order form a
    unitary circuit with operator $U^\dagger$ and the depth and gate count of $c$.
\end{theorem}
\begin{proof}\leanok
    The sequence of a concatenation of layers is the composition of the two
    sequences, whose Kraus operators multiply. The adjoint of a sequence of layers of
    single operators has the adjoint Kraus operator; a layer and its adjoint use the
    same cells, so they have the same costs.
\end{proof}

\begin{lemma}[One gate and one-site unitaries]
    \label{lem:unitary_circuit_gate}
    \lean{QuantumCircuit.UnitaryCircuit.ofLayer, QuantumCircuit.UnitaryCircuit.op_ofLayer,
        QuantumCircuit.UnitaryCircuit.depth_ofLayer,
        QuantumCircuit.UnitaryCircuit.gateCount_ofLayer,
        QuantumCircuit.Layer.forall_isUnitary_single,
        QuantumCircuit.Layer.forall_isUnitary_ofOnsite,
        QuantumCircuit.UnitaryCircuit.gate, QuantumCircuit.UnitaryCircuit.op_gate,
        QuantumCircuit.UnitaryCircuit.gateCount_gate,
        QuantumCircuit.UnitaryCircuit.depth_gate_le,
        QuantumCircuit.UnitaryCircuit.onsite, QuantumCircuit.UnitaryCircuit.op_onsite,
        QuantumCircuit.UnitaryCircuit.depth_onsite,
        QuantumCircuit.UnitaryCircuit.gateCount_onsite}
    \leanok
    \uses{def:unitary_circuit, lem:circuits_layer_single, def:sites_place}
    A unitary $V$ on the configurations of a cell $b$ is the unitary circuit of one
    layer with operator $V\otimes1$ placed on $C_b$, gate count $c_b$ and depth at
    most $1$. One-site unitaries $u_v$ at every site form a circuit of one layer with
    no cells, with operator $\bigotimes_vu_v$ and depth and gate count $0$.
\end{lemma}
\begin{proof}\leanok
    The layer of one gate has the placed matrix as Kraus operator and one cell; the
    layer of onsite families has the product of the onsite operators and no cell.
\end{proof}

\begin{theorem}[Light cone of a unitary circuit]
    \label{thm:unitary_circuit_light_cone}
    \lean{QuantumCircuit.UnitaryCircuit.conjTranspose_mul_mul_mem_supportedOperators_lightCone,
        QuantumCircuit.UnitaryCircuit.conjTranspose_mul_mul_mem_supportedOperators}
    \leanok
    \uses{def:unitary_circuit, lem:unitary_circuit_costs, thm:circuits_light_cone,
        def:sites_supported}
    If $A$ is supported on $R$, then $U^\dagger AU$ is supported on the backward light
    cone of $R$ through the circuit; when every cell has positive cost and the depth is
    at most $T$, on the $T$-fold neighbourhood $\mathcal N^T(R)$
    \cite[Theorem~1]{Malz2024Preparation}.
\end{theorem}
\begin{proof}\leanok
    The Heisenberg picture of the circuit is $A\mapsto U^\dagger AU$; the circuit is
    unbranched and trace-preserving, so the light cone of
    Chapter~\ref{ch:circuits} applies, and it lies in the $T$-fold neighbourhood
    when the cells have positive cost.
\end{proof}

\section{Depth and gate count}
\label{sec:unitary_classes}

\begin{definition}[Circuits of bounded depth and products of gates]
    \label{def:unitary_classes}
    \lean{QuantumCircuit.IsUnitaryCircuit, QuantumCircuit.IsGateProduct}
    \leanok
    \uses{def:unitary_circuit}
    A matrix $U$ on $\mathcal H_X$ is a unitary circuit of depth $T$ on $\Gamma$ if it
    is the operator of a unitary circuit of depth at most $T$
    \cite{Malz2024Preparation, Piroli2021LOCC}, and a product of gates of cost $K$ if
    it is the operator of a unitary circuit of gate count at most $K$. On the open
    chain these are the products of at most $K$ unitaries on neighbouring pairs of
    sites, with one-site unitaries free.
\end{definition}

\begin{theorem}[Closure of the two classes]
    \label{thm:unitary_classes_closure}
    \lean{QuantumCircuit.IsUnitaryCircuit.one, QuantumCircuit.IsUnitaryCircuit.mono,
        QuantumCircuit.IsUnitaryCircuit.mul, QuantumCircuit.IsUnitaryCircuit.conjTranspose,
        QuantumCircuit.IsUnitaryCircuit.mem_unitaryGroup, QuantumCircuit.IsUnitaryCircuit.pow,
        QuantumCircuit.IsUnitaryCircuit.prodOp, QuantumCircuit.IsUnitaryCircuit.place,
        QuantumCircuit.IsGateProduct.one, QuantumCircuit.IsGateProduct.mono,
        QuantumCircuit.IsGateProduct.mul, QuantumCircuit.IsGateProduct.conjTranspose,
        QuantumCircuit.IsGateProduct.mem_unitaryGroup, QuantumCircuit.IsGateProduct.pow,
        QuantumCircuit.IsGateProduct.list_prod, QuantumCircuit.IsGateProduct.prodOp,
        QuantumCircuit.IsGateProduct.place,
        QuantumCircuit.IsGateProduct.of_mem_supportedOperators,
        QuantumCircuit.IsGateProduct.isUnitaryCircuit}
    \leanok
    \uses{def:unitary_classes, thm:unitary_circuit_composition, lem:unitary_circuit_gate,
        lem:unitary_circuit_costs, thm:sites_place_range}
    Both classes contain $1$ and are monotone in the bound. Depths and gate counts add
    under products, are kept by adjoints and are multiplied by $k$ under $k$-th
    powers; a product of $\ell$ products of $K$ gates is a product of $\ell K$ gates.
    Products of one-site unitaries cost nothing, and a unitary supported on a cell $b$
    is one gate of cost $c_b$ and depth at most $1$. Every member is unitary, and a
    product of gates of cost $K$ is a circuit of depth $K$.
\end{theorem}
\begin{proof}\leanok
    By Theorem~\ref{thm:unitary_circuit_composition} and
    Lemma~\ref{lem:unitary_circuit_gate}. A unitary supported on $C_b$ is $V\otimes1$ for
    a matrix $V$ on $C_b$ (when every basis set is nonempty; otherwise every matrix is
    zero), and $V$ is unitary because placement is injective and multiplicative. The
    depth of a circuit is at most its gate count.
\end{proof}

\begin{theorem}[Light cone of a circuit of depth $T$]
    \label{thm:unitary_classes_light_cone}
    \lean{QuantumCircuit.IsUnitaryCircuit.conjTranspose_mul_mul_mem_supportedOperators,
        QuantumCircuit.IsUnitaryCircuit.conjTranspose_mul_mul_mem_supportedOperators_ringBall}
    \leanok
    \uses{def:unitary_classes, thm:unitary_circuit_light_cone,
        lem:circuits_ring_neighbourhood}
    On a geometry whose cells have positive cost, if $U$ is a unitary circuit of depth
    $T$ and $A$ is supported on $R$, then $U^\dagger AU$ is supported on
    $\mathcal N^T(R)$; on the ring of $N$ sites, on the sites within ring distance $T$
    of $R$ \cite[Theorem~1]{Malz2024Preparation}.
\end{theorem}
\begin{proof}\leanok
    Theorem~\ref{thm:unitary_circuit_light_cone}; on the ring the $T$-fold
    neighbourhood is the ring ball of radius $T$.
\end{proof}

\section{Preparation in depth $T$}
\label{sec:unitary_preparation}

\begin{definition}[Prepared in depth $T$]
    \label{def:unitary_prepared}
    \lean{QuantumCircuit.IsPreparedInDepth, QuantumCircuit.IsPreparedInDepth.mono,
        QuantumCircuit.IsPreparedInDepth.mulVec, QuantumCircuit.IsPreparedInDepth.prod}
    \leanok
    \uses{def:unitary_classes}
    A vector $\psi\in\mathcal H_X$ is prepared in depth $T$ if $\psi=U\bigotimes_v v_v$
    for a unitary circuit $U$ of depth $T$ and vectors $v_v\in\C^{X_v}$
    \cite{Malz2024Preparation}. A product vector is prepared in depth $0$, and a
    circuit of depth $T'$ applied to a vector prepared in depth $T$ gives one prepared
    in depth $T+T'$.
\end{definition}

\begin{lemma}[Product vectors over disjoint regions]
    \label{lem:unitary_product_vector}
    \lean{QuantumCircuit.dotProduct_prod_mul_mul}
    \leanok
    \uses{def:sites_supported, cor:sites_product_induction}
    For a product vector $\pi=\bigotimes_vv_v$ and operators $A$, $B$ supported on
    disjoint regions,
    $\bra\pi AB\ket\pi\braket\pi\pi=\bra\pi A\ket\pi\bra\pi B\ket\pi$.
\end{lemma}
\begin{proof}\leanok
    Both sides are bilinear in $(A,B)$, so it suffices to take product operators
    $A=\bigotimes_vm_v$, $B=\bigotimes_vm'_v$ with $m_v=1$ outside the first region and
    $m'_v=1$ outside the second. Then $\bra\pi\bigotimes_vm_v\ket\pi=
        \prod_v\bra{v_v}m_v\ket{v_v}$, and at every site one of $m_v$, $m'_v$ is the identity,
    so the two sides agree site by site.
\end{proof}

\begin{theorem}[No connected correlations beyond the light cone]
    \label{thm:unitary_correlations}
    \lean{QuantumCircuit.IsPreparedInDepth.dotProduct_mul_mul,
        QuantumCircuit.IsPreparedInDepth.dotProduct_mul_eq,
        QuantumCircuit.IsPreparedInDepth.dotProduct_mul_mul_of_isRingSeparated,
        QuantumCircuit.IsPreparedInDepth.dotProduct_mul_eq_of_isRingSeparated}
    \leanok
    \uses{def:unitary_prepared, lem:unitary_product_vector,
        thm:unitary_classes_light_cone, lem:sites_ring_separated}
    On a geometry whose cells have positive cost, let $\psi$ be prepared in depth $T$
    and let $A$, $B$ be supported on regions $R$, $R'$ with disjoint $T$-fold
    neighbourhoods. Then
    \begin{align}
        \bra\psi AB\ket\psi\braket\psi\psi=\bra\psi A\ket\psi\bra\psi B\ket\psi,
        \notag
    \end{align}
    and $\bra\psi AB\ket\psi=\bra\psi A\ket\psi\bra\psi B\ket\psi$ when $\psi$ is a unit
    vector. On the ring this holds for regions at ring distance larger than $2T$
    \cite[Supplemental Material, proof of Theorem~1]{Malz2024Preparation}.
\end{theorem}
\begin{proof}\leanok
    Write $\psi=U\pi$ with $\pi$ a product vector. Then
    $\bra\psi M\ket\psi=\bra\pi U^\dagger MU\ket\pi$ and
    $U^\dagger ABU=(U^\dagger AU)(U^\dagger BU)$. By
    Theorem~\ref{thm:unitary_classes_light_cone} the two factors are supported on the
    disjoint neighbourhoods, and Lemma~\ref{lem:unitary_product_vector} applies. On the
    ring, regions at ring distance larger than $2T$ have disjoint balls of radius $T$.
\end{proof}

\section{Substituting circuits for gates}
\label{sec:unitary_substitution}

\begin{lemma}[A unitary layer gate by gate]
    \label{lem:unitary_layer_factorization}
    \lean{QuantumCircuit.Layer.onsiteOp,
        QuantumCircuit.Layer.kraus_eq_prod_place_mul_prodOp,
        QuantumCircuit.Layer.onsiteOp_mem_unitaryGroup,
        QuantumCircuit.Layer.ev_gate_mem_unitaryGroup}
    \leanok
    \uses{def:circuits_layer, def:sites_place, thm:circuits_of_cells}
    Every Kraus operator of a layer is the product, over the cells in any order, of the
    matrices of the gates placed on their cells, times the tensor product of the onsite
    operators, with the identity at the sites covered by cells. For a unitary layer every
    factor is unitary.
\end{lemma}
\begin{proof}\leanok
    The product of the placed gates on disjoint cells is the Kraus operator of the layer of
    these gates with the identity elsewhere; multiplying by the product operator puts the
    onsite entries at the free sites, and the entries agree with those of the layer.
\end{proof}

\begin{theorem}[Substitution of gates]
    \label{thm:unitary_gate_substitution}
    \lean{QuantumCircuit.IsGateProduct.list_prod_map,
        QuantumCircuit.Layer.isGateProduct_kraus,
        QuantumCircuit.UnitaryCircuit.isGateProduct_op,
        QuantumCircuit.IsGateProduct.substitute}
    \leanok
    \uses{lem:unitary_layer_factorization, thm:unitary_classes_closure}
    Let $\Gamma$ and $\Gamma'$ be geometries on the same sites. If every unitary placed on a
    cell $b$ of $\Gamma$ is a product of gates of $\Gamma'$ of cost $f(b)$, the operator of a
    unitary circuit on $\Gamma$ is a product of gates of $\Gamma'$ of cost the sum of $f$ over
    the cells it uses. In particular, if $f(b)\le Cc_b$, a product of gates of cost $K$ on
    $\Gamma$ is a product of gates of cost $CK$ on $\Gamma'$. The depth admits no such bound:
    the circuits replacing gates on disjoint cells may use overlapping cells of $\Gamma'$.
\end{theorem}
\begin{proof}\leanok
    Write every layer as in Lemma~\ref{lem:unitary_layer_factorization}; the placed gates are
    replaced by their products on $\Gamma'$, the onsite product costs nothing, and costs add
    under products.
\end{proof}

\section{Swap routing on the open chain}
\label{sec:unitary_swap_routing}

On the open chain of $n$ sites with the same basis $\alpha$ at every site, the swap
of two sites is the operator of the transposition of the sites.

\begin{lemma}[Factoring a permutation into transpositions]
    \label{lem:unitary_swap_factors}
    \lean{Equiv.Perm.swapFactorsAux_length_le, Equiv.Perm.swapFactors_length_le_card}
    \leanok
    The transposition factorization of a permutation of a finite linearly ordered set uses
    at most one transposition per element.
\end{lemma}
\begin{proof}\leanok
    The factorization processes the elements in order and uses one transposition for each
    element not yet fixed.
\end{proof}

\begin{theorem}[Swap routing]
    \label{thm:unitary_swap_routing}
    \lean{QuantumCircuit.isGateProduct_sitePermOp_swap_of_succ,
        QuantumCircuit.isGateProduct_of_mem_supportedOperators_pair_of_lt,
        QuantumCircuit.isGateProduct_of_mem_supportedOperators_pair_dist,
        QuantumCircuit.isGateProduct_of_mem_supportedOperators_pair,
        QuantumCircuit.isGateProduct_sitePermOp_swap,
        QuantumCircuit.isGateProduct_sitePermOp,
        QuantumCircuit.IsGateProduct.sitePermOp_mul_mul_conjTranspose,
        QuantumCircuit.IsGateProduct.of_sitesAndPairs,
        QuantumCircuit.IsGateProduct.place_window_two,
        QuantumCircuit.IsGateProduct.of_window_two}
    \leanok
    \uses{def:unitary_graph_geometry, thm:unitary_classes_closure,
        thm:unitary_gate_substitution, lem:unitary_swap_factors,
        thm:sites-local-permutation}
    On the open chain, the swap of two neighbouring sites is one gate. A unitary supported
    on two sites $i\ne j$ is a product of at most $2|i-j|-1\le2n$ neighbouring gates, and so
    is the swap of $i$ and $j$; the operator of a permutation of the sites costs at most
    $2n^2$, and conjugation by it adds at most $4n^2$ to a gate count. A product of gates of
    cost $K$ on the geometry of all sites and pairs of sites, with pairs of cost at least
    $1$, is a product of at most $2nK$ neighbouring gates. On a chain of at least two sites a
    product of gates on windows of two sites is a product of the same number of neighbouring
    gates. This is the use of swap gates in \cite{Malz2024Preparation} (paragraph ``The
    sequential-RG circuit'': the inputs ``are separated by $O(q)$ sites, which requires one to
    implement SWAP gates'').
\end{theorem}
\begin{proof}\leanok
    A neighbouring swap is supported on its pair. For $i<j$, conjugating a unitary supported
    on $\{i,j-1\}$ by the neighbouring swap of $j-1$ and $j$ gives one supported on
    $\{i,j\}$, and induction on $j-i$ gives the count. A permutation is a product of at most
    $n$ transpositions by Lemma~\ref{lem:unitary_swap_factors}. The last two statements are
    Theorem~\ref{thm:unitary_gate_substitution}; one-site gates are free, and a window of two
    sites lies in a neighbouring pair.
\end{proof}

\section{Controlled operations and universality}
\label{sec:unitary_universality}

\begin{definition}[Controlled operations on sites]
    \label{def:unitary_ctrl_op}
    \lean{QuantumCircuit.ctrlProj, QuantumCircuit.ctrlOp}
    \leanok
    \uses{def:gate_ctrl_elem, def:sites_place_at, lem:sites_place_at_product}
    For a set $R$ of sites, a configuration $c$ and a site $t$, let $p_{R,c}$ be the
    projection onto the configurations that agree with $c$ on $R$, and let $u_t$ be a matrix
    $u$ on $\C^{X_t}$ placed at $t$. The controlled operation is
    $C_{R,c,t}(u)=p_{R,c}u_t+(1-p_{R,c})$: it applies $u$ at $t$ when the sites of $R$ are in
    the configuration $c$ \cite{Malz2024Preparation}.
\end{definition}

\begin{lemma}[Algebra of controlled operations]
    \label{lem:unitary_ctrl_op_algebra}
    \lean{QuantumCircuit.isIdempotentElem_ctrlProj, QuantumCircuit.commute_ctrlProj,
        QuantumCircuit.ctrlProj_insert, QuantumCircuit.ctrlProj_empty,
        QuantumCircuit.ctrlProj_conjTranspose, QuantumCircuit.ctrlProj_mem_supportedOperators,
        QuantumCircuit.commute_ctrlProj_placeAt, QuantumCircuit.ctrlOp_empty,
        QuantumCircuit.ctrlOp_mul, QuantumCircuit.ctrlOp_one,
        QuantumCircuit.ctrlOp_conjTranspose, QuantumCircuit.ctrlOp_mem_unitaryGroup,
        QuantumCircuit.ctrlOp_mem_supportedOperators,
        QuantumCircuit.ctrlOp_insert_commutator, QuantumCircuit.twoLevel_update_eq_ctrlOp}
    \leanok
    \uses{def:unitary_ctrl_op, thm:gate_ctrl_elem_algebra, thm:gate_ctrl_elem_commutator,
        def:gate_two_level}
    For $t\notin R$, the projections $p_{R,c}$ are commuting self-adjoint idempotents,
    $p_{R\cup\{s\},c}=p_{R,c}p_{\{s\},c}$, and $p_{R,c}$ commutes with $u_t$. Hence
    $C_{R,c,t}(u)C_{R,c,t}(w)=C_{R,c,t}(uw)$, $C_{R,c,t}(u)^\dagger=C_{R,c,t}(u^\dagger)$,
    $C_{R,c,t}(u)$ is unitary for unitary $u$ and is supported on $R\cup\{t\}$, and for
    $s\ne t$ and unitaries $v$, $w$,
    \begin{align}
        C_{R\cup\{s\},c,t}(vwv^\dagger w^\dagger)
        =C_{R,c,t}(v)\,C_{\{s\},c,t}(w)\,C_{R,c,t}(v^\dagger)\,C_{\{s\},c,t}(w^\dagger).
        \notag
    \end{align}
    The two-level operator between configurations $a$ and $a'$ that differ only at $t$ is
    $C_{V\setminus\{t\},a,t}$ of the two-level operator between the levels $a_t$ and $a'_t$.
\end{lemma}
\begin{proof}\leanok
    The projections are diagonal with entries in $\{0,1\}$. The commutator identity is that of
    controlled elements of a ring applied to the idempotents $p_{R,c}$ and $p_{\{s\},c}$. The
    two-level operator acts as the identity on the configurations differing from $a$ away
    from $t$, which is the controlled form.
\end{proof}

\begin{theorem}[Universality relative to the cost of two-site unitaries]
    \label{thm:unitary_universality}
    \lean{QuantumCircuit.quarterTwo_isSpecialTwo, QuantumCircuit.twoLevel_quarterTwo_mul,
        QuantumCircuit.twoLevel_conj_quarterTwo,
        QuantumCircuit.IsGateProduct.place_singleton, QuantumCircuit.IsGateProduct.smul_one,
        QuantumCircuit.IsGateProduct.of_mem_supportedOperators_singleton,
        QuantumCircuit.IsGateProduct.of_mem_supportedOperators_card_le_two,
        QuantumCircuit.isGateProduct_ctrlOp_twoLevel, QuantumCircuit.isGateProduct_twoLevel,
        QuantumCircuit.unitaryGateCount, QuantumCircuit.isGateProduct_unitaryGateCount,
        QuantumCircuit.exists_isGateProduct, QuantumCircuit.unitaryGateCount_le,
        QuantumCircuit.isGateProduct_sitesAndPairs_of_mem_supportedOperators_pair,
        QuantumCircuit.isGateProduct_sitesAndPairs}
    \leanok
    \uses{lem:unitary_ctrl_op_algebra, thm:gate_givens, thm:gate_special_two_commutator,
        thm:unitary_classes_closure, def:circuits_local_gate}
    Let $V$ be a nonempty finite set of $n$ sites and $m=\prod_v|X_v|$, and suppose that every
    unitary supported on two distinct sites is a product of gates of $\Gamma$ of cost $k$.
    One-site unitaries and global phases cost nothing. A controlled two-level rotation or
    phase with $j$ controls costs at most $(2^{j+1}-2)k$, a two-level rotation or phase
    between distinct configurations at most $(2n-1)(2^n-2)k$, and every unitary on
    $\mathcal H_X$ at most $3m^2(2n-1)(2^n-2)k\le6m^4k$ when every $|X_v|\ge2$. On the geometry
    of all sites and pairs of sites the hypothesis holds with $k$ the cost of a pair. This is
    the step ``can be further expressed with a low-depth circuit of local gates'' of
    \cite[caption of Fig.~1]{Malz2024Preparation}.
\end{theorem}
\begin{proof}\leanok
    A controlled operation with no control is a one-site unitary. Every rotation or phase of
    the special form is a commutator $g_1g_2g_1^\dagger g_2^\dagger$ with $g_1$ of the same
    form, so Lemma~\ref{lem:unitary_ctrl_op_algebra} removes one control at the cost of two
    controlled operations with one control, which are two-site unitaries. A two-level
    operator between configurations at Hamming distance $h+1$ is conjugated by a quarter
    turn on a pair at distance $1$ into one between configurations at distance $h$; with
    $n-1$ controls this gives the second count. A unitary is a phase times a unitary of
    determinant $1$, and the Givens decomposition writes the latter as at most $3m^2$
    two-level rotations and phases.
\end{proof}

\begin{theorem}[Universality of neighbouring gates on the open chain]
    \label{thm:unitary_chain_universality}
    \lean{QuantumCircuit.isGateProduct_pathGraph_unitaryGateCount,
        QuantumCircuit.isGateProduct_pathGraph_le_card_pow,
        QuantumCircuit.exists_isGateProduct_pathGraph_apply_embedding_eq,
        Matrix.exists_mem_unitaryGroup_apply_embedding_eq}
    \leanok
    \uses{thm:unitary_universality, thm:unitary_swap_routing}
    On the open chain of $n\ge1$ sites with the same basis at every site, every unitary is an
    exact product, global phase included and with no auxiliary sites, of at most
    $2n\cdot3m^2(2n-1)(2^n-2)$ neighbouring gates, and of at most $12m^5$ when the basis has
    at least two elements, $m$ the dimension of the Hilbert space. An isometry prescribed on a
    set of input columns extends to a unitary that is such a product.
\end{theorem}
\begin{proof}\leanok
    Theorem~\ref{thm:unitary_swap_routing} gives the hypothesis of
    Theorem~\ref{thm:unitary_universality} with $k=2n$, and $n<2^n\le m$. An isometry extends to
    a unitary with the prescribed columns.
\end{proof}


\section{Circuits inside a region and in parallel}
\label{sec:unitary_region}

\begin{definition}[A circuit inside a region]
    \label{def:unitary_inside}
    \lean{QuantumCircuit.Layer.IsIn, QuantumCircuit.UnitaryCircuit.IsIn,
        QuantumCircuit.IsUnitaryCircuitOn, KrausFamily.IsUnitary.op_eq_of_eq}
    \leanok
    \uses{def:unitary_circuit, def:unitary_classes}
    A layer acts inside a region $R$ when every cell it uses lies in $R$ and every onsite
    family at a site outside $R$ is the identity; a unitary circuit acts inside $R$ when all
    its layers do. A matrix $U$ is a unitary circuit of depth $T$ inside $R$ if it is the
    operator of a unitary circuit inside $R$ of depth at most $T$
    \cite{Malz2024Preparation}.
\end{definition}

\begin{lemma}[Circuits inside a region]
    \label{lem:unitary_inside}
    \lean{QuantumCircuit.Layer.IsIn.mono, QuantumCircuit.Layer.IsIn.not_mem_biUnion,
        QuantumCircuit.Layer.kraus_mem_supportedOperators, QuantumCircuit.Layer.isIn_single,
        QuantumCircuit.Layer.isIn_ofOnsite, QuantumCircuit.Layer.IsIn.adjoint,
        QuantumCircuit.UnitaryCircuit.tail, QuantumCircuit.UnitaryCircuit.op_eq_tail_op_mul,
        QuantumCircuit.UnitaryCircuit.op_eq_one, QuantumCircuit.UnitaryCircuit.IsIn.mono,
        QuantumCircuit.UnitaryCircuit.IsIn.tail, QuantumCircuit.UnitaryCircuit.isIn_id,
        QuantumCircuit.UnitaryCircuit.IsIn.comp, QuantumCircuit.UnitaryCircuit.IsIn.adjoint,
        QuantumCircuit.UnitaryCircuit.isIn_gate, QuantumCircuit.UnitaryCircuit.isIn_onsite,
        QuantumCircuit.UnitaryCircuit.op_mem_supportedOperators,
        QuantumCircuit.IsUnitaryCircuitOn.isUnitaryCircuit,
        QuantumCircuit.IsUnitaryCircuitOn.mem_unitaryGroup,
        QuantumCircuit.IsUnitaryCircuitOn.mem_supportedOperators,
        QuantumCircuit.IsUnitaryCircuitOn.mono_set, QuantumCircuit.IsUnitaryCircuitOn.mono,
        QuantumCircuit.IsUnitaryCircuitOn.one, QuantumCircuit.IsUnitaryCircuitOn.mul,
        QuantumCircuit.IsUnitaryCircuitOn.conjTranspose,
        QuantumCircuit.IsUnitaryCircuitOn.place,
        QuantumCircuit.IsUnitaryCircuitOn.of_mem_supportedOperators,
        QuantumCircuit.IsUnitaryCircuitOn.prodOp}
    \leanok
    \uses{def:unitary_inside, thm:unitary_classes_closure, def:sites_supported}
    The operator of a circuit inside $R$ is supported on $R$. Circuits inside $R$ are closed
    under composition and adjoints and stay inside every larger region; a gate on a cell in
    $R$, and one-site unitaries equal to the identity outside $R$, are circuits inside $R$.
\end{lemma}
\begin{proof}\leanok
    The Kraus operator of a layer inside $R$ is a product of operators supported on its cells
    and onsite operators that are the identity outside $R$.
\end{proof}

\begin{theorem}[Circuits in parallel]
    \label{thm:unitary_parallel}
    \lean{QuantumCircuit.Layer.par, QuantumCircuit.Layer.disjoint_cells,
        QuantumCircuit.Layer.forall_isUnitary_par, QuantumCircuit.Layer.isIn_par,
        QuantumCircuit.Layer.prod_free, QuantumCircuit.Layer.kraus_par,
        QuantumCircuit.Layer.gateCount_par, QuantumCircuit.UnitaryCircuit.parLayers,
        QuantumCircuit.UnitaryCircuit.forall_parLayers,
        QuantumCircuit.UnitaryCircuit.sum_depth_parLayers_le,
        QuantumCircuit.UnitaryCircuit.sum_gateCount_parLayers,
        QuantumCircuit.UnitaryCircuit.op_parLayers, QuantumCircuit.UnitaryCircuit.par,
        QuantumCircuit.UnitaryCircuit.isIn_par, QuantumCircuit.UnitaryCircuit.depth_par_le,
        QuantumCircuit.UnitaryCircuit.gateCount_par, QuantumCircuit.UnitaryCircuit.op_par,
        QuantumCircuit.IsUnitaryCircuitOn.par,
        QuantumCircuit.IsUnitaryCircuitOn.finset_noncommProd,
        QuantumCircuit.IsUnitaryCircuitOn.list_prod}
    \leanok
    \uses{lem:unitary_inside, lem:sites_supported_commute}
    On a geometry whose cells are nonempty, unitary circuits of depths $T_1$, $T_2$ inside
    disjoint regions $R_1$, $R_2$ give a unitary circuit inside $R_1\cup R_2$ with operator
    $U_1U_2$, depth at most $\max(T_1,T_2)$ and the sum of the gate counts. A finite family of
    circuits of depth $T$ inside pairwise disjoint regions has product a circuit of depth $T$
    inside their union. These are the parallel layers of \cite{Malz2024Preparation}
    (eq.~(10) and the paragraph ``The sequential-RG circuit'').
\end{theorem}
\begin{proof}\leanok
    Layers of the two circuits are merged greedily: a leading layer of depth $0$ is kept
    alone, and two leading layers of depth $1$ become one layer with the cells of both, which
    are disjoint since nonempty cells lie in disjoint regions. Depth then counts the merged
    pairs once, so the depth of the merged circuit is at most the larger depth. The operators
    of the two circuits commute, being supported on disjoint regions, so the operator of the
    merge is $U_1U_2$.
\end{proof}

\section{Placing and relabelling circuits}
\label{sec:unitary_placement}

\begin{theorem}[A circuit on a region placed in the system]
    \label{thm:unitary_place}
    \lean{QuantumCircuit.UnitaryCircuit.op_eq_of_forall₂, QuantumCircuit.Geometry.restrictHom,
        QuantumCircuit.Layer.place, QuantumCircuit.Layer.forall_isUnitary_place,
        QuantumCircuit.Layer.isIn_place, QuantumCircuit.Layer.gateCount_place,
        QuantumCircuit.Layer.depth_place, QuantumCircuit.Layer.kraus_place,
        QuantumCircuit.UnitaryCircuit.place, QuantumCircuit.UnitaryCircuit.op_place,
        QuantumCircuit.UnitaryCircuit.isIn_place, QuantumCircuit.UnitaryCircuit.depth_place,
        QuantumCircuit.UnitaryCircuit.gateCount_place, QuantumCircuit.IsGateProduct.of_restrict,
        QuantumCircuit.IsUnitaryCircuit.isUnitaryCircuitOn_of_restrict,
        QuantumCircuit.IsGateProduct.isUnitaryCircuitOn_of_restrict}
    \leanok
    \uses{def:unitary_inside, def:circuits_geometry, def:sites_place}
    A unitary circuit on the cells of $\Gamma$ inside a region $D$, on the system of the sites
    of $D$, read on all the sites with the identity outside $D$, is a unitary circuit inside
    $D$ of the same depth and gate count with operator $V\otimes1$ placed on $D$. In
    particular a product of gates of cost $K$ on the cells inside $D$ is, placed, a circuit of
    depth $K$ inside $D$ and a product of gates of cost $K$.
\end{theorem}
\begin{proof}\leanok
    Layer by layer, the cells and the onsite families at sites of $D$ are kept, and the
    identity is added at the sites outside $D$; the Kraus operator is the placement of that of
    the layer, and placement is multiplicative.
\end{proof}

\begin{theorem}[Relabelling sites]
    \label{thm:unitary_relabel}
    \lean{QuantumCircuit.Layer.relabel, QuantumCircuit.Layer.forall_isUnitary_relabel,
        QuantumCircuit.Layer.gateCount_relabel, QuantumCircuit.Layer.depth_relabel,
        QuantumCircuit.Layer.kraus_relabel, QuantumCircuit.UnitaryCircuit.relabel,
        QuantumCircuit.UnitaryCircuit.op_relabel, QuantumCircuit.UnitaryCircuit.depth_relabel,
        QuantumCircuit.UnitaryCircuit.gateCount_relabel, QuantumCircuit.IsGateProduct.relabel,
        QuantumCircuit.IsUnitaryCircuit.relabel, QuantumCircuit.Geometry.Hom.toRestrict,
        QuantumCircuit.Geometry.exists_hom_graph, QuantumCircuit.Geometry.graph_cell_nonempty,
        QuantumCircuit.Geometry.ring_cell_nonempty,
        QuantumCircuit.IsGateProduct.isUnitaryCircuitOn_of_hom,
        QuantumCircuit.IsGateProduct.isUnitaryCircuitOn_graph,
        QuantumCircuit.IsGateProduct.isUnitaryCircuitOn_ring,
        QuantumCircuit.IsGateProduct.place_pathGraph}
    \leanok
    \uses{thm:unitary_place, def:circuits_geometry_hom, lem:sites-regional-relabel,
        def:unitary_graph_geometry, def:circuits_ring_geometry}
    A bijection of sites carrying the cells of one geometry to cells of another of the same
    cost carries unitary circuits to unitary circuits with the relabelled operator and the
    same depth and gate count. Hence a product of $K$ gates on neighbouring sites of a chain of
    $m$ sites, placed on $m$ consecutive sites of a ring or of a longer chain, is a circuit of
    depth $K$ inside the placed sites and a product of $K$ neighbouring gates there
    \cite{Malz2024Preparation}.
\end{theorem}
\begin{proof}\leanok
    The layers are moved cell by cell along the bijection and the operators are relabelled
    along the induced bijection of configurations. Consecutive sites go to cells of the ring
    or of the chain, which gives a morphism of geometries into the cells inside the placed
    sites, and Theorem~\ref{thm:unitary_place} applies.
\end{proof}

\begin{theorem}[Gates on windows through neighbouring gates]
    \label{thm:unitary_window}
    \lean{QuantumCircuit.isGateProduct_pathGraph_of_mem_supportedOperators_interval,
        QuantumCircuit.IsGateProduct.window_of_mem_supportedOperators,
        QuantumCircuit.IsGateProduct.of_window}
    \leanok
    \uses{thm:unitary_relabel, thm:unitary_chain_universality, thm:unitary_gate_substitution}
    On the open chain with the same basis $\alpha$, $|\alpha|\ge2$, at every site, a unitary
    supported on a window of $m\ge1$ consecutive sites is one gate on the window and a product
    of at most $12(|\alpha|^m)^5$ neighbouring gates. A product of gates of cost $K$ on
    windows of $m$ sites is a product of at most $12(|\alpha|^m)^5K$ neighbouring gates, a
    factor independent of the length of the chain \cite[caption of Fig.~1]{Malz2024Preparation}.
\end{theorem}
\begin{proof}\leanok
    The window is a chain of at most $m$ sites placed on consecutive sites; universality on it
    (Theorem~\ref{thm:unitary_chain_universality}) and Theorem~\ref{thm:unitary_relabel} give
    the first count, and Theorem~\ref{thm:unitary_gate_substitution} the second.
\end{proof}

\section{Clean workspaces and reflections}
\label{sec:unitary_clean_workspace}

\begin{definition}[Clean implementation]
    \label{def:unitary_clean_implementation}
    \lean{QuantumCircuit.IsCleanImplementation}
    \leanok
    \uses{def:gate_basis_inclusion}
    For an inclusion $J$ of a logical space into a larger space with its workspace in an
    initial state, an operator $U$ cleanly implements $Z$ when $UJ=JZ$: it acts as $Z$ on
    every logical input, including inputs entangled with a reference, and returns the
    workspace to its initial state.
\end{definition}

\begin{lemma}[Closure of clean implementations]
    \label{lem:unitary_clean_closure}
    \lean{QuantumCircuit.IsCleanImplementation.one, QuantumCircuit.IsCleanImplementation.mul,
        QuantumCircuit.IsCleanImplementation.pow,
        QuantumCircuit.IsCleanImplementation.conjTranspose,
        QuantumCircuit.IsCleanImplementation.inv, QuantumCircuit.IsCleanImplementation.mulVec,
        QuantumCircuit.IsCleanImplementation.logical_mem_unitaryGroup,
        QuantumCircuit.IsCleanImplementation.conjTranspose_of_isometry,
        QuantumCircuit.IsCleanImplementation.kronecker_one}
    \leanok
    \uses{def:unitary_clean_implementation}
    Clean implementations are closed under products and powers, under adjoints and inverses
    when both operators are unitary, and under tensoring with the identity of a reference. A
    unitary that cleanly implements $Z$ along an isometric inclusion makes $Z$ unitary, so
    its adjoint cleanly implements $Z^\dagger$.
\end{lemma}
\begin{proof}\leanok
    $UVJ=UJW=JZW$. From $UJ=JZ$ and $ZZ^\dagger=1$, $U^\dagger J=U^\dagger JZZ^\dagger=
        U^\dagger UJZ^\dagger=JZ^\dagger$. If $J^\dagger J=1$, then
    $Z^\dagger Z=(JZ)^\dagger JZ=(UJ)^\dagger UJ=1$.
\end{proof}

\begin{theorem}[Clean implementations by products of gates]
    \label{thm:unitary_clean_existence}
    \lean{QuantumCircuit.exists_isGateProduct_mul_basisInclusion_eq,
        QuantumCircuit.exists_isGateProduct_isCleanImplementation,
        QuantumCircuit.workspaceEmbedding, QuantumCircuit.workspaceEmbedding_apply,
        QuantumCircuit.basisInclusion_workspaceEmbedding_mul_apply,
        QuantumCircuit.workspaceFixedSitePerm, QuantumCircuit.workspaceFixedSitePerm_castAdd,
        QuantumCircuit.workspaceFixedSitePerm_natAdd,
        QuantumCircuit.isCleanImplementation_sitePermOp_workspaceFixedSitePerm,
        QuantumCircuit.IsCleanImplementation.conjugate_workspaceFixedSitePerm}
    \leanok
    \uses{lem:unitary_clean_closure, thm:unitary_chain_universality,
        thm:sites-local-permutation}
    If every unitary is a product of gates of cost $K$, every isometry on prescribed input
    columns and every logical unitary, along a basis inclusion, is cleanly implemented by a
    product of gates of cost $K$. A permutation of the logical sites that fixes an appended
    workspace cleanly implements the permutation of the logical register, and so does its
    conjugate of a clean implementation, so one workspace can be reused by many calls.
\end{theorem}
\begin{proof}\leanok
    The isometry extends to a unitary with the prescribed columns, which is a product of
    gates by hypothesis. The site permutation leaves the workspace sites unchanged, and the
    identity is checked entry by entry; conjugation uses Lemma~\ref{lem:unitary_clean_closure}.
\end{proof}

\begin{theorem}[Reflections about implemented images]
    \label{thm:unitary_image_reflection}
    \lean{Matrix.subspaceReflection_image, QuantumCircuit.IsGateProduct.subspaceReflection_image,
        QuantumCircuit.IsGateProduct.postselectionAmplificationStep}
    \leanok
    \uses{thm:unitary_classes_closure, def:gate_subspace_amplification,
        def:gate_uniform_postselection}
    For a unitary $U$, the reflection about the image of $UE$ is $U$ times the reflection
    about the image of $E$ times $U^\dagger$. If $U$ is a product of gates of cost $K$ and the
    reflection about the image of $E$ one of cost $L$, the reflection about the image of $UE$
    costs $2K+L$; one round of amplification by two reflections of costs $K$ and $L$, with its
    global phase $-1$, costs $K+L$ on a nonempty set of sites.
\end{theorem}
\begin{proof}\leanok
    $UEE^\dagger U^\dagger$ is the projection onto the image of $UE$, and conjugation by $U$
    commutes with $1-2P$. The global phase is a one-site phase.
\end{proof}

\begin{theorem}[Conversion of encoders through a common seed]
    \label{thm:unitary_encoder_conversion}
    \lean{QuantumCircuit.IsUnitaryCircuit.norm_encoder_conversion_le}
    \leanok
    \uses{thm:unitary_classes_closure}
    If $U_A$ and $U_B$ are unitary circuits of depths $T_A$ and $T_B$ and $J$ is any seed,
    then $U_BU_A^\dagger$ is a unitary circuit of depth $T_A+T_B$ and
    $\|U_BU_A^\dagger F-G\|\le\|F-U_AJ\|+\|G-U_BJ\|$ in the operator norm
    \cite{Malz2024Preparation}.
\end{theorem}
\begin{proof}\leanok
    $U_BU_A^\dagger F-G=U_BU_A^\dagger(F-U_AJ)+(U_BJ-G)$, and $U_BU_A^\dagger$ has norm at
    most $1$.
\end{proof}

\begin{theorem}[Preparation from a product basis vector]
    \label{thm:unitary_preparation_basis}
    \lean{Matrix.exists_mem_unitaryGroup_eq_smul_mulVec_single,
        QuantumCircuit.exists_prod_eq_smul_prodOp_mulVec,
        QuantumCircuit.IsPreparedInDepth.exists_eq_smul_mulVec_prod_single}
    \leanok
    \uses{def:unitary_prepared, thm:unitary_classes_closure}
    Every nonzero vector of $\C^\iota$ is $c\,u\ket{i_0}$ with $c\ne0$ and $u$ unitary, every
    nonzero product vector is $c\bigotimes_vu_v\ket{x_0}$, and every nonzero vector prepared
    in depth $T$ is $c\,U\ket{x_0}$ for a unitary circuit $U$ of depth $T$ and any fixed
    product basis vector $\ket{x_0}$: the one-site unitaries cost no depth.
\end{theorem}
\begin{proof}\leanok
    Extend $x/\|x\|$ to a unitary with this column at $i_0$; take the product over the sites;
    and compose the circuit with the one-site unitaries, of depth $0$.
\end{proof}
